\documentclass[12pt,letterpaper]{article}
\usepackage[utf8]{inputenc}
\usepackage[spanish]{babel}
\usepackage{amsmath}
\usepackage{amsfonts}
\usepackage{amssymb}
\usepackage[margin=2.5cm]{geometry}

\begin{document}

\begin{center}
    \Large \textbf{Evaluación de Examen 2: Unidad II} \\
    \large Física para Ciencias de la Salud (005-1713) \\
    \vspace{0.5cm}
    \textbf{Estudiante:} Mariana Salazar \quad \textbf{C.I.:} 33.570.679
    \vspace{0.3cm} \\
    \large \textbf{Calificación Final: 10/10 puntos}
\end{center}

\vspace{0.5cm}

\section*{Análisis de Resultados}

\subsection*{1. Cinemática (Aceleración media)}
\textbf{Procedimiento y resultado:} Extrae ordenadamente los datos y plantea la ecuación de aceleración media ($a = \frac{V_f - V_i}{t}$). La sustitución y cálculo de $0,6 / 0,15$ arroja el resultado analítico exacto de $4 \text{ m/s}^2$ con un impecable manejo de unidades en su desarrollo. \\
\textbf{Veredicto:} \textbf{Correcto} (2/2 pts).

\subsection*{2. Dinámica (Leyes de Newton)}
\textbf{Procedimiento y resultado:} Muestra un despeje correcto a partir de la Segunda Ley de Newton para hallar la aceleración ($a = \frac{F}{m}$). Efectúa la sustitución de $150 \text{ N}$ entre $25 \text{ kg}$, logrando el resultado perfecto de $6 \text{ m/s}^2$. \\
\textbf{Veredicto:} \textbf{Correcto} (2/2 pts).

\subsection*{3. Trabajo Mecánico}
\textbf{Procedimiento y resultado:} Extrae los datos (identificando acertadamente que el ángulo es $0^\circ$) e identifica la fórmula general del trabajo. Efectúa correctamente la sustitución ($400 \text{ N} \cdot 0,8 \text{ m} \cdot \cos(0^\circ)$). El cálculo aritmético da exactamente $320 \text{ J}$. \\
\textbf{Veredicto:} \textbf{Correcto} (2/2 pts).

\subsection*{4. Energía Potencial}
\textbf{Procedimiento y resultado:} Extrae rigurosamente los datos del problema (masa, altura y constante gravitacional). Plantea el producto de las tres variables según la ecuación de la energía potencial gravitatoria ($E_p = m \cdot g \cdot h$) obteniendo el resultado matemático exacto de $17,64 \text{ J}$. \\
\textbf{Veredicto:} \textbf{Correcto} (2/2 pts).

\subsection*{5. Potencia Mecánica}
\textbf{Procedimiento y resultado:} Extrae los datos y relaciona correctamente el trabajo efectuado y el tiempo planteando la ecuación respectiva ($P = \frac{W}{t}$). El desarrollo de la división de $75 \text{ J} / 60 \text{ s}$ arroja la cifra correcta y verificada de $1,25 \text{ Watts}$. \\
\textbf{Veredicto:} \textbf{Correcto} (2/2 pts).

\vspace{0.5cm}
\section*{Comentarios Generales y Sugerencias}
¡Felicidades por obtener la nota máxima! El examen demuestra un entendimiento sólido de los conceptos de la Unidad II y un gran nivel de estructuración y limpieza en la resolución de problemas.
\begin{itemize}
    \item El formato implementado en cada hoja (separación explícita de \textit{Datos}, \textit{Fórmula} y \textit{Resultado} encuadrado) facilita enormemente la corrección y el seguimiento metodológico de los cálculos. Es una excelente práctica clínica y de laboratorio.
    \item Los cálculos aritméticos, magnitudes y análisis dimensionales fueron manejados de forma muy acertada, integrando las unidades incluso durante las sustituciones intermedias.
    \item Como un pequeñísimo apunte formal de cara a futuros documentos técnicos: la convención estricta del Sistema Internacional indica que el símbolo de \textit{Watts} se escribe en mayúscula ($\text{W}$) o, de escribir la palabra completa, esta no se pluraliza ($\text{Watt}$).
    \item ¡Sigue con este mismo entusiasmo y nivel de excelencia en las próximas unidades!
\end{itemize}

\end{document}